\section{Introduction}\label{intro}

The core of this project was to understand the key principles of deep learning and to apply them in a practical environment. This was achieved through two challenges. The first was to classify the digits 0~-~9 in the MNIST dataset, while the second challenge was to classify nine different tissue patterns in the PathMNIST dataset. A crucial first step in developing an effective classification model was to thoroughly investigate and understand the dataset at hand. Therefore, our investigation began with an introduction to the well-known datasets MNIST \citep{deng2012mnist} and MedMNIST \citep{medmnistv1}, which served as building blocks for our study.

The project was structured to take account of the different characteristics of the individual datasets. We started with MNIST, which was chosen due to its wide distribution and the numerous tutorials available, which simplified our entry into the field of deep learning. With this dataset, we took on the challenge of developing a rudimentary Convolutional Neural Network (SimpleCNN) that was intentionally designed with a limited number of layers. The initial aim of this challenge was not to achieve peak performance, but rather to gain practical experience and understand the basics of the architecture of neural networks and their ability to distinguish between different digits.

As our knowledge increased, we shifted our focus to the more challenging MedMNIST dataset, focusing particularly on the PathMNIST subset. This phase formed the core of our project, in which we focused intensively on experimenting with different pre-trained models. Our investigations extended to testing a wide range of hyperparameters as well as implementing different strategies for data preprocessing and augmentation. 

The report is structured as follows: First, the datasets used~-~MNIST and MedMNIST~-~are briefly introduced. We then review some theoretical foundations that are important for understanding the methods used in each challenge, such as the dropout layer or ReLU.\@ This is followed by an explanation of different architectures of deep learning models, starting from our SimpleCNN to more advanced models such as AlexNet, ResNet and Xception. We then present the results obtained with each model for the respective challenges. In the discussion that follows, we critically evaluate our results for each model. Finally, we briefly reflect on the lessons we have learned from the project.

\subsection{MNIST}\label{mnist}
The ``modified National Institute of Standards and Technology'' dataset comprises a collection of 70,000 handwritten digits carefully divided into a training set of 60,000 images and a test set of 10,000 images. Each digit is represented in a grayscale image of 28$\times$28 pixels, and offers a wide range of styles and shapes. This dataset is widely recognized for its simplicity and effectiveness in benchmarking classification algorithms, making it an ideal starting point for those new to deep learning. Some examples of the dataset can be seeen in Figure~\ref{fig:MNIST}

\begin{figure}
    \centering
    \includegraphics[width=0.3\textwidth]{figures/MNIST.png}
    \caption{Example greyscale images of the MNIST dataset~-~CHANGE IMAGE~-~RENDER OWN!!!}\label{fig:MNIST}
\end{figure}

Chosen for our initial challenge, MNIST provided a fundamental platform to explore neural network basics and experiment with simple model architectures. It allowed us to grasp the essentials of model training, and hyperparameter testing.


\subsection{MedMNIST}\label{medmnist}

MedMNIST, a more specialized and challenging dataset than MNIST, is tailored for medical image classification tasks. It extends the concept of handwritten digit classification to a diverse range of medical imaging modalities, including dermatology or radiology. Unlike MNIST's uniform format, MedMNIST encompasses 12 subsets for 2D and 6 subsets for 3D data. For our project, we focused on the PathMNIST \citep{kather2018, kather2019} subset, which includes ``100.000 non-overlapping image patches from hematoxylin and eosin stained histological images, and a test dataset [$\ldots$] of 7.180 image patches from a different clinical center'' \citep{medmnistv1}. The images could be classified into nine different types of tissues.

Initially, the images were of high resolution (3 $\times$ 224 $\times$ 224 pixels), but the authors of MedMNIST resized them to 3 $\times$ 28 $\times$ 28 pixels. The 100.000 training images were then divided into training and validation sets in a 9:1 ratio. Some examples of the various images can be seen in Figure~\ref{fig:MedMNIST}.

\begin{figure}
    \centering
    \includegraphics[width=0.3\textwidth]{figures/MNIST.png}
    \caption{Example images of the MedMNIST dataset.~-~CHANGE IMAGE~-~RENDER OWN!!!}\label{fig:MedMNIST}
\end{figure}


The PathMNIST subset of MedMNIST2D provides a unique challenge by introducing the complexity of medical image analysis. It requires the use of advanced deep learning techniques and models to accurately classify different types of tissue, making it an excellent progression from the simpler MNIST dataset.

